% =====================================================================
% 7 Plano de trabalho, riscos e contingências
% Fontes: docs/ars/19 §5 (matriz R1–R8) e §7; docs/ars/18;
% docs/legislacao/lai-*-minuta.md; docs/ars/17 §5.
% =====================================================================

\chapter{Plano de trabalho e riscos}
\label{cap:plano}

\section{Atividades}
\label{sec:atividades}

O trabalho restante se divide em cinco frentes.

\subsection{Medida do tratamento}

A restrição ativa do Ensaio 1 é saber onde e quando havia pulverização aérea
antes da lei. Três atividades a enfrentam.

\begin{alineas}
  \item Protocolar os pedidos de acesso à informação já minutados: ao MAPA,
    pelos relatórios mensais dos operadores aeroagrícolas com atuação no Ceará
    (o texto da Instrução Normativa nº 2/2008 será conferido antes do
    protocolo); à ADAGRI, pelo receituário agronômico com modalidade de
    aplicação; e à SEMACE, pelo cadastro de empresas prestadoras de serviço.
  \item Consultar, na Justiça Federal no Ceará, a ação civil pública de 2011
    contra empresas de fruticultura da Chapada do Apodi, para saber se houve
    liminar ou sentença que tenha interrompido a pulverização aérea antes de
    2019.
  \item Conferir o texto integral da Lei Municipal nº 1.511/2010 de Limoeiro do
    Norte.
\end{alineas}

\subsection{Consolidação dos dados e da inferência}

\begin{alineas}
  \item Estender o painel até 19 de dezembro de 2024 e registrar a data de
    extração de cada base;
  \item executar com dados reais a inferência sobre o nível que respeita os 15
    controles, o perfil por faixa de dose e o intervalo por inversão do teste
    de permutação, já implementados;
  \item reestimar a dose com a janela de 2010 a 2014 e medir a correlação de
    postos entre as duas janelas.
\end{alineas}

\subsection{Desenho com mais poder}

O projeto avaliou oito rotas de continuação. A mais promissora com os dados
existentes é um desenho de calendário: comparar, nos municípios onde havia
aeronave, a diferença de desfecho entre gestações expostas aos meses de pico de
aplicação e gestações expostas à entressafra, antes e depois da lei, contra a
mesma diferença nos demais municípios. A hipótese de identificação passa a ser
a estabilidade desse hiato sazonal, e não o paralelismo entre municípios com e
sem banana. Se os relatórios do MAPA forem obtidos, os municípios tratados e os
meses de pico deixam de ser supostos e passam a ser medidos.

\decisao{7}{adotar o desenho de calendário como etapa seguinte do Ensaio 1.
Ele muda a estratégia de identificação e exige plano de análise próprio,
fechado antes de cruzar exposição e desfecho.} A sequência prevista é: calcular
o efeito mínimo detectável do desenho de calendário apenas com dados anteriores
à lei, com o código já escrito; decidir com o orientador se o poder justifica a
adoção; fechar e versionar o plano de análise; e só então estimar.

\subsection{Ensaio 2}

Leitura do acórdão da ADI 6137 e dos autos das ADI 5553 e 7755; levantamento da
dispersão do custo de conversão de aplicação aérea para terrestre; redação das
três partes descritas na Seção~\ref{sec:e2-enquadramento}.

\subsection{Redação e defesa}

Revisão do texto do Ensaio 1 conforme o parecer recebido em setembro de 2026;
redação do Ensaio 2; introdução e conclusão gerais; depósito e defesa.

\section{Riscos e contingências}
\label{sec:riscos}

O Quadro~\ref{qua:riscos} relaciona os riscos principais.

\begin{quadro}[htb]
\centering
\caption{Riscos e contingências}
\label{qua:riscos}
\small
\begin{tabular}{|p{4.4cm}|p{4.2cm}|p{5.5cm}|}
\hline
\textbf{Risco} & \textbf{Consequência} & \textbf{Contingência} \\
\hline
Os pedidos de acesso à informação são negados ou os registros não existem
  & O tratamento continua medido por \emph{proxy}
  & Desenho de calendário com os municípios do Censo de 2006 e meses de pico
    declarados como hipótese \\
\hline
A pulverização aérea cessou na Chapada em 2011
  & Não há tratamento a avaliar onde ele se concentrava
  & O Ensaio 1 é apresentado como diagnóstico do desenho e da medida, com o
    achado documental como resultado \\
\hline
O desenho de calendário não tem poder
  & O Ensaio 1 permanece inconclusivo quanto ao efeito
  & Reportar o efeito de forma reduzida com inferência adequada e o limite de
    poder como contribuição metodológica \\
\hline
A série até 2024 altera as estimativas
  & Números do Capítulo~\ref{cap:resultados} mudam
  & As duas versões são reportadas, com data de extração \\
\hline
Prazo
  & Atraso do Ensaio 2
  & O Ensaio 2 é analítico e não depende de dado novo; pode ser redigido em
    paralelo \\
\hline
\end{tabular}
\legend{Fonte: elaboração própria.}
\end{quadro}

Em qualquer cenário, o Ensaio 1 tem um resultado defensável: o diagnóstico de
que o desenho disponível mede um efeito de forma reduzida, com o poder e a
qualidade da medida do tratamento quantificados contra critério fixado
previamente. Os cenários favoráveis acrescentam a esse diagnóstico uma
estimativa com tratamento medido.

\section{Reprodutibilidade e ética}
\label{sec:reprodutibilidade}

O código do projeto é versionado, com execução ordenada por um arquivo de
orquestração, ambientes fixados para Python e R e suíte de testes automatizada.
A preparação dos dados é feita em Python e a estimação com tratamento contínuo
em R, com troca de arquivos entre as duas etapas. O plano de análise e seus
desvios são versionados com data.

Os dados são secundários, públicos e anonimizados na origem (DATASUS, IBGE,
FAO). Os microdados não são redistribuídos, e os resultados são reportados em
agregados municipais. \pendente{confirmar com o programa se o projeto, por
usar apenas dados públicos sem identificação, dispensa apreciação por comitê
de ética, nos termos da Resolução CNS nº 510/2016.}
